% Draft results paragraph — "Cross-loop representational dynamics".
% Forward-only P0 probe on the mid-trained base (osrt_v5_midtrain_final.pt).
% Numbers from research/probe_cross_loop_kv_{general,mixed}.json + the math run.
% Slot into paper.tex (analysis/results section). Not yet committed to paper.tex.

\subsection{Cross-loop representational dynamics}

We ask whether the recursive stack puts \emph{distinct} information in the
per-loop key/value latents, both to test the recursion-as-iterative-refinement
premise and to gate a candidate optimization: collapsing the $L{=}6$ per-loop KV
caches of each physical block into one (a potential $6\times$ decode-bandwidth
reduction stacking on KDV's $2\times$). The probe is forward-only --- a single
\texttt{use\_cache=True} pass over a held-out batch, no training --- and dumps,
per physical block, the cross-loop linear CKA of the cached latents
$\mathbf{c}_{\text{kv}}$, the per-token cosine, and the \emph{injected} relative
$L_2$ error a reuse scheme would impose.

\paragraph{The reuse hypothesis is falsified.} Forcing every loop to reuse
loop~0's latent injects a mean relative $L_2$ error of $\approx\!1.02\text{--}1.04$
--- the perturbation is as large as the signal itself --- and a two-group split
(cache 2 of 6) still costs $0.50\text{--}0.97$. Block~0 is the strongest evidence:
its post-loop-0 cosine to loop~0 is \emph{negative} ($\approx\!-0.4$) while CKA
stays structured ($0.52$), i.e.\ later loops occupy a related subspace under a
sign-flip, not noise. Cross-loop KV sharing is therefore rejected for this model;
a low-rank per-loop delta is the only variant left, and only if decode bandwidth
later becomes binding.

\paragraph{The recursion is a non-degenerate contracting iteration.} The same
data yields a positive structural result. Adjacent-loop similarity rises
monotonically along the chain, so the per-step representational move
$1-\mathrm{CKA}(k,k{+}1)$ \emph{shrinks} --- in block~0,
$\{0.30, 0.17, 0.07, 0.05, 0.05\}$ --- and the late loops converge to a near
fixed point ($\mathrm{CKA}(4,5)\!\approx\!0.95$). Loop~0 acts as an encoding pass
(the sign-flip above); loops $1{\to}5$ form a diminishing refinement trajectory.
Three independent signals captured on the \emph{same} forward agree: (i) the
KV-CKA contraction above; (ii) the residual update magnitude
$\lVert\Delta x\rVert/\lVert x\rVert$ is strongly front-loaded
($12.7 \to 0.9 \to \dots \to 0.5$ across loops in block~0); and (iii) the MoE
\emph{per-token} routing entropy rises from $0.91$ nats at loop~0 to
$\approx\!1.98$ ($\ln 8 = 2.08$) at later loops --- the router commits sharply in
the encoding pass and then refines diffusely, a genuine per-token spread rather
than mere aggregate load balance. The effect is carried by block~0; the deeper
physical blocks route broadly throughout and instead show a small terminal
uptick at the final loop, consistent with interior contraction followed by
output shaping.

\paragraph{Robustness.} The structure is stable across three disjoint probe
batches (math, general-domain, mixed): mean $\mathrm{CKA}(\cdot,\text{loop}_0) =
0.656 / 0.687 / 0.656$ and mean first-share $L_2 = 1.04 / 1.02 / 1.04$, and the
finding independently reproduced on a separately seeded run. The contraction is
thus not an artifact of the math-heavy default text.

\paragraph{Corollary: a narrow but safe inference knob.} The contraction
motivates the variable-loop knob (\S\ref{sec:variable-loops}), but representation
convergence is \emph{not} output equivalence; we measure the output shift
directly via token-level $\mathrm{KL}\!\left(P_{L=6}\,\|\,P_{L=k}\right)$ on the
same batches. Only the final loop is cheap to drop ($6\!\to\!5$: KL $0.04$,
$+6\%$ ppl), and the per-dropped-loop cost is \emph{superlinear} ($\Delta$KL
$=0.04, 0.06, 0.08, 0.21$) --- the exact inverse of the contraction curve: late
loops do the least representational work yet sit nearest the LM head, where the
map has highest gain. The usable knob is thus \emph{one} loop, not the 2--3 the
CKA plateau invites. Top-1 agreement ($0.85$ at $6\!\to\!5$) \emph{overstates}
the damage: the flipped tokens are low-confidence near-ties (full-model
confidence $0.18$ at flipped vs $0.46$ at unflipped positions), so the faithful
cost is the small KL/ppl move. Categorizing the flips by gold-token type sharpens
it further --- the last loop \emph{never} changes a mathematical operator or
logical connective and under-weights digits (type enrichment $0.0/0.0/0.3$ vs
$1.16$ for generic words), revising only uncertain function-word choices.
Reasoning-critical tokens are therefore settled by loop~5 and left untouched by
loop~6, making the drop safe precisely where reasoning lives --- though it also
means the last loop's marginal computation is lexical polish, not reasoning, so
we make no claim that depth is spent \emph{on} reasoning tokens. The per-loop
auxiliary LM heads (\S\ref{sec:aux-heads}), which train every depth to be
decodable, are the mechanism by which a representationally near-redundant late
loop still emits a coherent prediction. All thresholds are reported on the
current mid-trained base (general/mixed batches agree to within $\sim\!1\%$ ppl);
the droppable depth should be re-measured after SFT/RL, since a better-trained
model may reach its fixed point --- and lock its reasoning tokens --- at a
different loop. The accuracy check (does the reasoning-on$>$off delta survive
fewer loops?) is deferred to a checkpoint exhibiting such a delta, since a
compute saving that erodes it is rejected regardless of throughput.
